% Part: sets-functions-relations
% Chapter: functions
% Section: partial-functions

\documentclass[../../../include/open-logic-section]{subfiles}


\begin{document}

\olfileid[tr]{sfr}{fun}{par}

\olsection{Kısmi Fonksiyonlar}

\begin{explain}
Bazen fonksiyon tanımını, fonksiyonun çıktısının bütün olası girdiler için
tanımlı olmasını gerektirmeyecek biçimde gevşetmek yararlıdır. Bu tür
eşlemelere \emph{kısmi fonksiyonlar} denir.
\end{explain}

\begin{defn}
Bir \emph{kısmi fonksiyon} $f \colon A \pto B$, $A$'nın her elemanına
$B$'nin en fazla bir elemanını atayan bir eşlemedir. $f$, bir $x \in A$
elemanına $B$'nin bir elemanını atıyorsa $f(x)$'in \emph{tanımlı}, aksi
hâlde \emph{tanımsız} olduğunu söyleriz. $f(x)$ tanımlıysa
$f(x) \fdefined$, değilse $f(x) \fundefined$ yazarız. Bir $f$ kısmi
fonksiyonunun \emph{tanım kümesi}, $A$'nın fonksiyonun tanımlı olduğu alt
kümesidir; yani $\dom{f} = \Setabs{x \in A}{f(x) \fdefined}$.
\end{defn}

\begin{ex}
Her $f\colon A \to B$ fonksiyonu aynı zamanda bir kısmi fonksiyondur.
$A$'nın her yerinde tanımlı olan, yani şimdiye kadar yalnızca fonksiyon
dediğimiz kısmi fonksiyonlara \emph{tümel} fonksiyonlar da denir.
\end{ex}

\begin{ex}
$f(x)=1/x$ eşitliğiyle verilen $f \colon \Real \pto \Real$ kısmi fonksiyonu
$x=0$ için tanımsız, diğer her yerde tanımlıdır.
\end{ex}

\begin{prob}
$f\colon A \pto B$ verildiğinde $g\colon B \pto A$ kısmi fonksiyonunu şöyle
tanımlayın: Her $y \in B$ için $f(x)=y$ olacak biçimde biricik bir
$x \in A$ varsa $g(y)=x$; aksi hâlde $g(y) \fundefined$ olsun. $f$ bire bir ise
her $x \in \dom{f}$ için $g(f(x))=x$ ve her $y \in \ran{f}$ için
$f(g(y))=y$ olduğunu gösterin.
\end{prob}

\begin{defn}[Bir kısmi fonksiyonun grafiği]
$f\colon A \pto B$ bir kısmi fonksiyon olsun. $f$'nin \emph{grafiği},
\[
R_f = \Setabs{\tuple{x,y}}{f(x) = y}
\]
eşitliğiyle tanımlanan $R_f \subseteq A \times B$ bağıntısıdır.
\end{defn}

\begin{prop}
$R \subseteq A \times B$ bağıntısının, $Rxy$ ve $Rxy'$ olduğunda $y=y'$
olması özelliğini taşıdığını varsayalım. O hâlde $R$, şu şekilde tanımlanan
$f\colon A \pto B$ kısmi fonksiyonunun grafiğidir: $Rxy$ olacak biçimde bir
$y$ varsa $f(x)=y$; aksi hâlde $f(x) \fundefined$ olur. $R$ ayrıca
\emph{seri} ise, yani her $x \in A$ için $Rxy$ olacak biçimde bir $y \in B$
varsa, $f$ tümeldir.
\end{prop}

\begin{proof}
$Rxy$ olacak biçimde bir $y$ bulunduğunu varsayalım. $Rxy'$ olacak biçimde
başka bir $y'\neq y$ bulunsaydı $R$ üzerindeki koşul çiğnenirdi. Dolayısıyla
$Rxy$ olacak biçimde bir $y$ varsa bu $y$ biriciktir ve böylece $f$ iyi
tanımlıdır. Açıktır ki $R_f=R$'dir ve $R$ seri ise $f$ tümeldir.
\end{proof}

\end{document}
